% P10-Evaluationsfeinschliff, 16.09.2026: Aussagegrenzen und Anforderungsbilanz; keine neuen Messwerte.
% B-59, 15.09.2026: datierte Ledger-Nachauswertung, Stützung und Reichweite.
% Beleg: thesis-evidence/w2/ledger-repetitions-20260915/ (Originale unverändert).
% P10-4b, 14.09.2026: Begriffs-/Beleganschlüsse und abgegrenzte Prüfaussagen präzisiert.
% ============================================================
% 8.2 — kompakte Arbeitsfassung, Phase 02b, 10.09.2026.
% Grundlage: 7.8 und die befundbezogenen Grenzen aus 7.3–7.7.
% Hier Konsequenzen für die Schlussführung; keine zweite Aufzählung
% der Bewertungsrunden oder der technischen Fehlerklassen.
% ============================================================

\section{Grenzen und Übertragbarkeit}
\label{sec:diskussion-grenzen-uebertragbarkeit}

Die in Abschnitt~\ref{sec:eval-grenzen} dokumentierten
Untersuchungsbedingungen begrenzen vor allem die Verallgemeinerung der
Befunde. Zwei entwicklungsbekannte Quellenfälle und ein Hauptlauf je
Bedingung charakterisieren konkrete Ergebnisse. Die ergänzenden
Ledger-Ausgaben beschreiben beobachtete Schwankungen an einer dieser
Quellen; die geringe Laufzahl erlaubt keine belastbare Aussage über
die typische Leistung. Die rechnerische Reproduktion der Kennzahlen
bleibt von der Stabilität der Modellausgaben zu unterscheiden. Die KI-gestützte Bildung und
Bewertung des Referenzbestands sowie die fehlende unabhängige
menschliche Zweitannotation lassen zudem Unsicherheit in den
Inhaltsurteilen bestehen. Deshalb sind
die beobachteten Unterschiede als Ergebnisse dieser Konfigurationen und
Fälle zu lesen; eine allgemeine Rangfolge der Ansätze folgt daraus nicht.

Für das Gesamtsystem ist zwischen nachgewiesener Realisierbarkeit und
umfassender Absicherung zu unterscheiden. Die Belege stammen aus
unterschiedlichen Systemständen und ausgewählten Ausführungspfaden.
Sie bestätigen keinen durchgängig geprüften späteren Endstand. Auch
der Governance-Anspruch bleibt pfadbezogen: Der initiale Seed wird ohne
Einzelautorisierung gespeichert; die nachgelagerten
Fortschreibungspfade besitzen eigene Freigaben. Das umfassende Autorisierungsziel A6 wird damit nur teilweise
erreicht. Diese Einschränkung betrifft auch die Zielvorstellung
eines autorisierten Bestands in A5, A7 und A8. Ungeprüfte Fehlerpfade und fehlende
Messungen der menschlichen Bearbeitung begrenzen dagegen den
Untersuchungsumfang. Ob die bereitgestellten Prüfmöglichkeiten tatsächlich
weniger menschliche Arbeit oder bessere Endergebnisse bewirken, bleibt
offen. Der Tokenaufwand allein beantwortet diese Frage nicht.

Eine Übertragung erscheint insbesondere für Arbeitskontexte prüfenswert,
in denen wiederkehrende Kommunikation einen bestehenden fachlichen Bestand
verändert und Herkunft sowie Übernahmeentscheidungen nachvollziehbar
bleiben müssen. Zu prüfen wären dabei die geeignete Granularität der
Evidenz, die erforderlichen fachlichen Beziehungen und die Grenzen
automatisierter Entscheidungen. Die konkrete Ausgestaltung hängt von
Domäne, Eingaben und Verantwortlichkeiten ab. Die Arbeit liefert dafür
begründete Gestaltungsansätze und benannte Fehlermöglichkeiten; die
Übertragbarkeit ihrer Leistungswerte und die Zweckmäßigkeit derselben
Architektur in anderen Kontexten bedürfen eigener Untersuchungen.
